\documentclass[11pt,letterpaper,oneside]{report}
\usepackage[T1]{fontenc}
\usepackage[utf8]{inputenc}
\usepackage[english]{babel}
\usepackage{lmodern}
\usepackage{microtype}
\usepackage[margin=0.86in,headheight=16pt]{geometry}
\usepackage[table,dvipsnames]{xcolor}
\usepackage{graphicx}
\usepackage{booktabs,tabularx,longtable,array}
\usepackage{enumitem}
\usepackage{titlesec}
\usepackage{fancyhdr}
\usepackage{hyperref}
\usepackage{xurl}
\Urlmuskip=0mu plus 1mu\relax
\usepackage{tcolorbox}
\usepackage{listings}
\usepackage{tikz}
\usetikzlibrary{arrows.meta,positioning,shapes.geometric,fit,calc}
\usepackage{needspace}
\usepackage{amssymb}
\definecolor{ink}{HTML}{14253A}
\definecolor{navy}{HTML}{17385D}
\definecolor{blue}{HTML}{20739A}
\definecolor{pale}{HTML}{ECF4F8}
\definecolor{green}{HTML}{1D6C5C}
\definecolor{softgreen}{HTML}{EEF8F4}
\definecolor{border}{HTML}{D4DFE7}
\definecolor{grey}{HTML}{506174}
\hypersetup{colorlinks=true,linkcolor=navy,urlcolor=blue,pdfauthor={Micro-Frontend Engineering Documentation Suite}}
\titleformat{\chapter}[display]{\normalfont\sffamily\bfseries\color{ink}}{\Large\color{blue}\MakeUppercase{\chaptername} \thechapter}{8pt}{\Huge}
\titlespacing*{\chapter}{0pt}{-10pt}{24pt}
\titleformat{\section}{\Large\sffamily\bfseries\color{navy}}{\thesection}{0.7em}{}
\titleformat{\subsection}{\large\sffamily\bfseries\color{navy}}{\thesubsection}{0.6em}{}
\setlist[itemize]{itemsep=3pt,topsep=5pt,leftmargin=1.6em}
\setlist[enumerate]{itemsep=4pt,topsep=5pt,leftmargin=1.8em}
\pagestyle{fancy}\fancyhf{}
\fancyhead[L]{\footnotesize\sffamily MICRO-FRONTEND ENGINEERING SUITE}
\fancyhead[R]{\footnotesize\sffamily\docshort}
\fancyfoot[L]{\footnotesize\color{grey}Engineering edition | October 2026}
\fancyfoot[R]{\footnotesize\thepage}
\renewcommand{\headrulewidth}{0.35pt}
\renewcommand{\arraystretch}{1.2}
\setlength{\emergencystretch}{2em}
\lstdefinestyle{eng}{basicstyle=\ttfamily\footnotesize,keywordstyle=\color{navy}\bfseries,commentstyle=\color{grey},stringstyle=\color{green},breaklines=true,breakatwhitespace=false,columns=fullflexible,keepspaces=true,showstringspaces=false,frame=single,rulecolor=\color{border},backgroundcolor=\color{pale},xleftmargin=5pt,xrightmargin=5pt,aboveskip=10pt,belowskip=10pt,captionpos=b}
\lstdefinelanguage{JavaScript}{morekeywords={const,let,var,require,return,async,await,import,export,default,function,try,catch,new,true,false,undefined,module},sensitive=true,morecomment=[l]//,morestring=[b]',morestring=[b]"}
\lstset{style=eng}
\newtcolorbox{sourcebox}{colback=pale,colframe=border,arc=2pt,boxrule=0.6pt,left=8pt,right=8pt,top=7pt,bottom=7pt}
\newtcolorbox{engineerbox}{colback=softgreen,colframe=green!45,arc=2pt,boxrule=0.6pt,left=8pt,right=8pt,top=7pt,bottom=7pt}
\newcommand{\based}[1]{\begin{sourcebox}\textbf{BOOK-GROUNDED CONCEPT}\\#1\end{sourcebox}}
\newcommand{\added}[1]{\begin{engineerbox}\textbf{ENGINEERING EXTENSION}\\#1\end{engineerbox}}
\newcommand{\term}[1]{\textbf{\color{navy}#1}}
\newcommand{\criterion}[1]{\textbf{Acceptance:} #1}
\newcommand{\book}{\emph{Building Micro-Frontends}, 2nd ed.}
\newcommand{\refs}{\chapter*{References and Provenance}\addcontentsline{toc}{chapter}{References and Provenance}
\begin{enumerate}[leftmargin=*]
\item \textbf{Primary source:} Luca Mezzalira, \emph{Building Micro-Frontends: Distributed Systems for the Frontend}, 2nd ed., O'Reilly Media, published November 2025 (copyright 2026), ISBN 978-1-098-17078-3. Chapter and printed-page references are included throughout. This suite paraphrases and synthesizes concepts; it does not reproduce the book.
\item \textbf{Engineering extension:} Webpack Module Federation reference: \url{https://webpack.js.org/plugins/module-federation-plugin}.
\item \textbf{Engineering extension:} Next.js multi-zones: \url{https://nextjs.org/docs/app/guides/multi-zones}. Details are version-sensitive and must be checked against the deployed framework version.
\item \textbf{Engineering extension:} OWASP Application Security Verification Standard 5.0: \url{https://owasp.org/www-project-application-security-verification-standard/}.
\item \textbf{Engineering extension:} GitHub Actions secure use reference: \url{https://docs.github.com/en/actions/reference/security/secure-use}.
\end{enumerate}}

\titleclass{\chapter}{straight}\titleformat{\chapter}[hang]{\huge\sffamily\bfseries\color{navy}}{\thechapter}{.7em}{}\titlespacing*{\chapter}{0pt}{24pt}{13pt}
\newcommand{\docshort}{IMPLEMENTATION}
\begin{document}
\begin{titlepage}\thispagestyle{empty}
\vspace*{1.1in}{\Huge\sffamily\bfseries\color{navy} Implementation Handbook\par}
\vspace{18pt}{\Large\sffamily\color{blue} TypeScript contracts, composition approaches, shell design, routing, and deployment patterns\par}
\vspace{27pt}\noindent\color{blue}\rule{0.75\textwidth}{1.4pt}
\vspace{23pt}\par{\large\sffamily A companion to Luca Mezzalira's \emph{Building Micro-Frontends}, Second Edition}\par
\vfill\noindent\textbf{Document series:} Micro-Frontend Software Architecture and Engineering\\
\textbf{Version:} 1.0 \quad \textbf{Prepared:} October 2026\\
\textbf{Production format:} pdfLaTeX-compatible standalone source\\
\textbf{Scope:} Original synthesis plus explicitly identified engineering extensions
\end{titlepage}
\pagenumbering{roman}\tableofcontents\clearpage\pagenumbering{arabic}

\chapter{Implementation Baseline}
\based{The source offers a Module Federation client-side application example (Chapter 4), SSR and multi-zone approaches (Chapter 5), and backend integration patterns (Chapter 9). The code below is independently written illustrative scaffolding, not a verbatim reproduction or a production-ready repository.}
\section{Example domain boundaries}
Build a storefront with \term{catalog}, \term{account}, and \term{checkout}. Start with account and catalog as route-owning domains. Add checkout only after the first domain has passed all CI, auth, and rollback gates. The shell manages navigation and release resolution, not cart pricing or customer authorization.
\section{Repository blueprint}
\begin{lstlisting}[caption={Illustrative monorepo layout}]
workspace/
  apps/shell/                  # Navigation and remote loading
  apps/catalog/                # Catalog UI and domain contract
  apps/account/                # Account UI and routes
  packages/ui-tokens/          # Stable design tokens
  packages/domain-contracts/   # Versioned types and schemas
  platform/discovery/          # Environment/release mapping
  platform/observability/      # Tracing conventions
  .github/workflows/           # CI and deployment automation
\end{lstlisting}
Monorepo does not imply synchronous deployment: produce immutable independently promotable artifacts. Polyrepo can work as well when the platform provides cross-repository integration tests and version discovery.
\chapter{Runtime Contract and Safe Composition}
\section{Domain lifecycle contract}
\begin{lstlisting}[caption={TypeScript interface for a client-side domain}]
export type AppContext = Readonly<{
  route: string;
  correlationId: string;
  events: EventTarget;
}>;

export interface DomainModule {
  mount(root: HTMLElement, ctx: AppContext): Promise<void>;
  unmount(root: HTMLElement): Promise<void>;
}
\end{lstlisting}
Apply lifecycle teardown to event listeners, subscriptions, and timers. Define how a second navigation during an asynchronous \texttt{mount} is cancelled. The shell never passes bearer credentials in event detail, query parameters, or browser global variables.
\section{Versioned event contracts}
\begin{lstlisting}[caption={Typed application event with minimal payload}]
type ItemSelectedV1 = {
  type: 'catalog.item-selected.v1';
  itemId: string;
};

function publishItemSelected(bus: EventTarget, itemId: string) {
  bus.dispatchEvent(new CustomEvent<ItemSelectedV1>('domain:event', {
    detail: { type: 'catalog.item-selected.v1', itemId }
  }));
}
\end{lstlisting}
The receiver may act on the identifier but must fetch/authorize its own data. Browser events are not a security boundary; validate payload shape and treat untrusted event sources accordingly.
\chapter{Webpack 5 Module Federation}
\section{Remote configuration}
\begin{lstlisting}[caption={webpack.config.js for the catalog remote},language=JavaScript]
const { container: { ModuleFederationPlugin } } = require('webpack');
module.exports = {
  mode: 'production',
  output: { publicPath: 'auto', uniqueName: 'catalogApp' },
  plugins: [new ModuleFederationPlugin({
    name: 'catalog',
    filename: 'remoteEntry.js',
    exposes: { './App': './src/catalog-entry.ts' },
    shared: { react: { singleton: true, requiredVersion: '^19.0.0' },
              'react-dom': { singleton: true, requiredVersion: '^19.0.0' } }
  })]
};
\end{lstlisting}
The dependency version is \emph{illustrative}; match it to the actual lockfile and test incompatible ranges. On the host, a remote mapping can be statically configured per environment or resolved via an authenticated/trusted deployment registry. Never load a remote URL merely because untrusted query input names it.
\section{Host contract and failure boundary}
\begin{lstlisting}[caption={Conceptual host import with fallback},language=JavaScript]
async function loadCatalog() {
  try {
    // Host federation config declares: catalog -> approved remote URL
    const remote = await import('catalog/App');
    return remote;
  } catch (error) {
    reportFailure('catalog-unavailable', error);
    return import('./fallback/catalog-unavailable');
  }
}
\end{lstlisting}
Add a bounded timeout at the loader/network boundary in a real host. A failed remote must not make the global navigation unreachable. Test cache consistency between \texttt{remoteEntry.js} and the referenced immutable asset chunks.
\chapter{Server-side and Route-based Alternatives}
\section{Multi-zone router model}
\begin{lstlisting}[caption={Conceptual reverse proxy route ownership}]
/account/*   -> account-zone (SSR)
/catalog/*   -> catalog-zone (SSR)
/*           -> default-shell-zone
/static/*    -> versioned CDN asset namespace
\end{lstlisting}
Route precedence, asset namespaces, and cookie boundaries must be verified for each framework version. For Next.js, follow the matching version's multi-zone documentation and test \texttt{assetPrefix}, \texttt{basePath}, rewrites, direct links, and hard versus soft navigations before standardizing a template.
\section{HTML fragment/server composition}
A composition gateway can fetch independently owned fragments, assemble a document, and stream the response. Set per-fragment timeouts, escape data at the correct output context, enforce a trusted HTML producer policy, and make caching explicit. Prefer server-side authorization for personalized fragments; shared CDN caching of user-specific HTML is a critical leakage risk.
\chapter{Discovery, Promotion, and Compatibility}
\based{Chapter 7 describes frontend discovery as a release-management and composition technique, distinct from domain feature flags.}
\begin{lstlisting}[caption={Illustrative discovery manifest (not the book schema)}]
{
  "schemaVersion": 1,
  "environment": "production",
  "domains": {
    "catalog": {
      "release": "2026.10.17+sha.example",
      "entry": "https://assets.example.test/catalog/abc123/remoteEntry.js",
      "integrityMetadataId": "manifest-record-123",
      "minimumShellContract": 2
    }
  }
}
\end{lstlisting}
This document invents no signing scheme: integrity metadata only helps if a trusted process validates it. Implement registry write authorization, authenticated updates, staged rollout rules, monotonic audit logs, cache expiration, and emergency reversion. Use a compatibility matrix before switching the active entry.
\chapter{Backend-for-Frontend and Identity}
\section{Contract boundaries}
The BFF terminates authenticated sessions or validates tokens at a trusted backend layer; it may aggregate domain APIs and transform payloads for its client. Return only fields the calling principal may see. Any domain that receives a record key must independently enforce authorization to that record.
\begin{lstlisting}[caption={Illustrative HTTP flow}]
Browser -> GET /api/account/profile (session cookie)
BFF     -> validate session and anti-CSRF policy as appropriate
BFF     -> request profile with server-held service credentials
Service -> authorize subject for profile record
BFF     -> return minimized JSON with no-store cache directive
\end{lstlisting}
Set Secure, HttpOnly, appropriate SameSite, path/domain scoping on cookies. Reconcile login/logout behavior across all zones. Never accept a shell-only route guard as an authorization control.
\chapter{Design System, A11y, and CSS Isolation}
\section{Shared visual grammar}
Share versioned design tokens, stable typography/spacing, and accessibility conventions. Avoid introducing a framework-coupled mega-library that forces all teams into coordinated dependency upgrades. Horizontal composition should test CSS collisions, stacking contexts, focus order, and interactions between nested scroll containers.
\section{End-to-end keyboard test plan}
\begin{enumerate}\item Reach shell navigation and each domain with keyboard alone.
\item Verify focus transitions on route change, dialog open/close, and error fallback.
\item Verify accessible names, headings, landmarks, and announcement of failures.
\item Verify login/logout does not strand focus in an unloaded domain.\end{enumerate}
\chapter{Caching, Performance, and Operations}
Cache immutable hashed JS assets long-term; give the mutable manifest/remote pointers intentional short TTL or explicit invalidation; validate asset origin availability before promotion. Segment personalized HTML and API responses by correct user/session rules, or mark them non-cacheable. Measure transfer sizes, long tasks, LCP, INP, CLS, error rate, and route-level reliability on representative client devices.
\section{Minimum telemetry fields}
Capture \texttt{route}, \texttt{domain}, \texttt{release}, \texttt{shellVersion}, \texttt{traceId/correlationId}, \texttt{environment}, and relevant HTTP timings, with privacy controls on IDs and events. Attribute third-party errors separately where possible.
\chapter{Implementation Acceptance Checklist}
\begin{itemize}[label=$\square$]\item Domain locally runnable without all remotes.
\item Contract test verifies module or route identity and schema.
\item Runtime dependency policy and CSS isolation verified.
\item Login, logout, direct deep link, session timeout, and forbidden API tested.
\item Fallback tested with inaccessible remote entry and missing chunk.
\item Deployment registry disallows unauthorized origin changes.
\item Asset cache and rollback behavior tested under two versions.
\item Browser keyboard, accessibility, and performance audits passed.
\item Every deployment links to immutable artifact and change provenance.\end{itemize}

\refs
\end{document}
